\documentclass[10pt]{article}
\usepackage[a4paper,margin=2cm]{geometry}
\usepackage[T1]{fontenc}
\usepackage[utf8]{inputenc}
\usepackage{tgpagella,mathpazo,amsmath,booktabs,longtable,array}
\usepackage[scaled=0.95]{helvet}
\usepackage{microtype,titlesec}
\usepackage[dvipsnames,table]{xcolor}
\usepackage{tikz}
\usetikzlibrary{arrows.meta}
\usepackage{enumitem}
\usepackage[numbers,sort&compress]{natbib}
\usepackage[hidelinks]{hyperref}
% Shared manuscript and TikZ visual language for the OPERA paper.
% Keep colours and Helvetica-like figure typography synchronized with
% opera/visualization/style.py (Arial/Liberation Sans on render hosts).

\definecolor{OperaIndigo}{HTML}{2D3A8C}
\definecolor{OperaNavy}{HTML}{1A237E}
\definecolor{OperaPeriwinkle}{HTML}{5C7AEA}
\definecolor{OperaSlate}{HTML}{737D8C}
\definecolor{OperaTeal}{HTML}{408579}
\definecolor{OperaGreen}{HTML}{1B7A4A}
\definecolor{OperaSienna}{HTML}{B5451B}
\definecolor{OperaTerracotta}{HTML}{C16A58}
\definecolor{OperaOchre}{HTML}{B18B4D}
\definecolor{OperaLavender}{HTML}{8974AC}
\definecolor{OperaInk}{HTML}{17212B}
\definecolor{OperaInkSecondary}{HTML}{485463}
\definecolor{OperaInkMuted}{HTML}{657180}
\definecolor{OperaBorder}{HTML}{CBD2DA}
\definecolor{OperaGrid}{HTML}{E8ECF1}
\definecolor{OperaPaper}{HTML}{F7F8FA}

\tikzset{
  opera/.style={font=\sffamily\small, text=OperaInk},
  panel/.style={draw=none, fill=white, inner sep=0pt},
  softpanel/.style={
    rounded corners=0.8mm,
    draw=OperaGrid,
    line width=0.45pt,
    fill=OperaPaper,
    inner sep=2.4mm,
  },
  stage/.style={
    rounded corners=0.7mm,
    draw=OperaBorder,
    line width=0.6pt,
    fill=white,
    align=center,
    minimum height=8mm,
    inner xsep=3mm,
    inner ysep=1.8mm,
  },
  stageblue/.style={stage, draw=OperaIndigo!55, fill=OperaIndigo!7},
  stageteal/.style={stage, draw=OperaTeal!65, fill=OperaTeal!8},
  stagewarm/.style={stage, draw=OperaSienna!60, fill=OperaSienna!7},
  latent/.style={
    circle,
    draw=OperaIndigo!65,
    fill=OperaIndigo!9,
    line width=0.6pt,
    minimum size=6mm,
    inner sep=0pt,
    font=\sffamily\scriptsize\bfseries,
  },
  token/.style={
    rounded corners=0.7mm,
    draw=OperaBorder,
    fill=white,
    minimum width=6.2mm,
    minimum height=4.8mm,
    inner sep=0.5mm,
    font=\sffamily\tiny,
    align=center,
  },
  outcome/.style={
    rounded corners=1mm,
    draw=OperaTerracotta!70,
    fill=OperaTerracotta!9,
    minimum height=5.8mm,
    inner xsep=1.8mm,
    font=\sffamily\scriptsize,
    align=center,
  },
  flow/.style={
    -{Latex[length=2mm,width=1.2mm]},
    line width=0.75pt,
    draw=OperaInkSecondary,
  },
  flowblue/.style={flow, draw=OperaIndigo},
  flowteal/.style={flow, draw=OperaTeal},
  flowwarm/.style={flow, draw=OperaSienna},
  relation/.style={
    -{Latex[length=1.8mm,width=1.1mm]},
    line width=0.65pt,
    draw=OperaLavender,
    densely dashed,
  },
  captionnote/.style={
    font=\sffamily\fontsize{6.6}{7.5}\selectfont,
    text=OperaInkMuted,
    align=center,
  },
  labelsmall/.style={
    font=\sffamily\fontsize{7}{8}\selectfont,
    text=OperaInkSecondary,
    align=center,
  },
  heading/.style={
    font=\sffamily\fontsize{9}{10}\selectfont\bfseries,
    text=OperaInk,
    align=left,
  },
  panelletter/.style={font=\sffamily\bfseries\normalsize, text=OperaInk},
  divider/.style={draw=OperaBorder, line width=0.45pt},
  eyebrow/.style={font=\sffamily\fontsize{6.5}{7.5}\selectfont\bfseries,
    text=OperaIndigo, align=left},
}


\raggedbottom
\renewcommand{\arraystretch}{1.22}
\setlength{\LTpre}{8pt}
\setlength{\LTpost}{10pt}
\titleformat{\section}{\sffamily\Large\bfseries\color{OperaIndigo}}{\thesection}{0.6em}{}
\titleformat{\subsection}{\sffamily\normalsize\bfseries\color{OperaInk}}{\thesubsection}{0.6em}{}
\titlespacing*{\section}{0pt}{18pt}{8pt}
\titlespacing*{\subsection}{0pt}{12pt}{5pt}
\setlength{\parindent}{0pt}
\setlength{\emergencystretch}{2.5em}
\setlength{\parskip}{4pt plus 1pt minus 1pt}
\setlist{nosep,leftmargin=*}
\renewcommand{\thesection}{S\arabic{section}}
\renewcommand{\thetable}{S\arabic{table}}
\renewcommand{\tablename}{Supplementary Table}

% Column layout for the outcome tables (Supplementary Tables S5--S11).
\newcolumntype{L}[1]{>{\raggedright\arraybackslash}p{#1}}
\newcommand{\otablebegin}{\begin{longtable}{L{2.75cm}L{5.35cm}L{1.55cm}L{2.45cm}L{2.25cm}L{1.0cm}}}
\newcommand{\otableheadrow}{\toprule
\rowcolor{OperaIndigo!8}
Outcome & Definition and threshold & Source & Baseline requirement & Ascertainment and censoring & Class\\
\midrule}
\newcommand{\otablehead}{%
\otableheadrow\endfirsthead
\otableheadrow\endhead
\bottomrule\endfoot}

\begin{document}
\ifdefined\embeddedSupplement\pagestyle{empty}\fi
\section*{Supplementary Methods: cohort, prediction time and outcome definitions}

This supplement describes how the study cohort, the index date (start of
first-line treatment), the temporal split, the disease-course endpoints and the
69 treatment-related and acute-care outcomes were defined in the Danish
Lymphoid Cancer Research (DALY-CARE) data resource~\citep{brieghel2025dalycare}.
Reported cohort counts describe the current release and await final supervisor
sign-off. Implementation traceability and unresolved clinical/source decisions
are recorded in the companion traceability and open-question documents.

\paragraph{Conventions.} Dates are calendar dates without time of day unless
stated otherwise. The \emph{index date} is the start of the reconstructed
first-line treatment episode (Section~\ref{sec:population}). \emph{End of
follow-up} is the earlier of death and last follow-up as recorded in the
DALY-CARE patient table. The \emph{administrative censoring date} is the
earlier of end of follow-up and the data cutoff of 19 March 2026. Unless stated
otherwise, an outcome event is the \emph{first} qualifying date $d$ with
$\text{index date} < d \le \text{censoring date}$; events on the index date
itself were never counted. For each outcome, patients for whom the outcome
could not be ascertained were excluded from that outcome's risk set (masked
supervision); exclusion was never interpreted as absence of the event.

% ===========================================================================
\paragraph{How to read this supplement.}
Sections~\ref{sec:sources}--\ref{sec:split} describe the data and who enters the
study, including the reconstructed treatment index and disease assignment.
Section~\ref{sec:2l} explains disease-course outcomes, particularly second-line
treatment. Section~\ref{sec:outcomes} gives the clinical outcome dictionary;
Section~\ref{sec:elig} explains who can contribute a label and when observation
ends. The final sections distinguish panel selection and reporting groups from
the clinical definitions and state their limitations.

\begin{center}
\fcolorbox{OperaBorder}{OperaPaper}{\begin{minipage}{0.92\linewidth}
\textbf{Three questions for every outcome}\par
\textbf{What counts?} The qualifying event and its threshold or code rule.\par
\textbf{Who is eligible?} Required baseline information and source coverage.\par
\textbf{How long is it observable?} Ascertainment, censoring and competing death.\par
An unavailable label is not a negative outcome. Clinical rationale and known
limitations accompany the implementation; unresolved source decisions remain
listed in the companion open-question document.
\end{minipage}}
\end{center}
\tableofcontents
\clearpage

\section{Data sources and linkage}\label{sec:sources}

All sources were accessed within the DALY-CARE secure environment and linked
at patient level by the DALY-CARE pseudonymised patient identifier.
Supplementary Table~\ref{tab:sources} lists the sources and their use.

{\footnotesize
\setlength{\tabcolsep}{3pt}
\begin{longtable}{L{3.3cm}L{4.3cm}L{4.0cm}L{4.4cm}}
\caption{Data sources.}\label{tab:sources}\\
\toprule
Source & Content used & Used for & Coverage and date handling\\
\midrule\endfirsthead
\toprule
Source & Content used & Used for & Coverage and date handling\\
\midrule\endhead
DALY-CARE patient table & Vital status, date of death or end of follow-up, sex, date of birth & Death, end of follow-up, demographics & End of follow-up is the earlier of death and last follow-up.\\
DALY-CARE haematological diagnoses & Harmonised haematological diagnoses with date, reliability rank and data source & Diagnosis anchor, disease classification, new haematological malignancy & Diagnoses from the Danish clinical quality registries (RKKP) were assigned the highest reliability rank.\\
Danish National Patient Registry: diagnoses & All ICD-10 diagnoses (SKS diagnosis codes) & New malignancy, cardiovascular, infection and fracture outcomes & LPR2 until 2019 and LPR3 thereafter; no era restriction.\\
DALY-CARE ICD-10 disease dictionaries & ICD-10 code patterns per haematological disease; SKS radiotherapy codes & Disease classification; radiotherapy as first-line treatment & Supplemented locally (Section~\ref{sec:disease}).\\
Harmonised antineoplastic treatment data (version of 10 September 2026) & Drug code (ATC or SKS), drug text, indication code and text, source, start date & First- and second-line treatment episodes & Treatment-plan records were excluded; records from the LYFO, DaMyDa and DCLLR quality registries were regarded as reliable.\\
Curated CLL and MM treatment-line data (CLL version 4, MM version 5) & Dated treatment lines & First-line anchoring in CLL/MM; second-line events in CLL/SLL and MM & Assumed complete for CLL/SLL and MM patients.\\
Danish Lymphoma Registry (LYFO) & Date of second-line treatment & Second-line treatment in lymphoma & Data from the CLL (DCLLR) and myeloma (DaMyDa) registries were not used for second-line dating.\\
Laboratory results & Analytes identified by NPU codes, cleaned with the DALY-CARE standard procedure & Laboratory outcomes and laboratory ascertainment & Complete capture assumed from 1 January 2017 (Section~\ref{sec:elig}).\\
DALY-CARE acute kidney injury data & Start dates of acute kidney injury episodes derived from creatinine by a separate DALY-CARE algorithm & Acute kidney injury outcome & Laboratory ascertainment rule.\\
Danish National Patient Registry: admissions, contacts and procedures (LPR2, LPR3) & Admissions and outpatient/emergency contacts (LPR2); physical contacts (LPR3); surgical and other procedures (SKS) & Procedure outcomes; acute hospitalisation & Procedures dated by the procedure start date, not the contact start.\\
Hospital Medicines Register (SMR) & In-hospital medicine administrations with ATC code and route & Intravenous antimicrobial-use outcome & Reporting starts 15 May 2018; outcome restricted to index dates on or after this date.\\
Danish National Prescription Registry & Dispensed prescriptions with ATC code and dispensing date & Prescription antimicrobial-use outcome & No era restriction.\\
International Prognostic Index (IPI) & IPI at diagnosis & Descriptive population characteristics only & Not an outcome or model input.\\
\bottomrule
\end{longtable}}

\paragraph{Linkage of procedures.} Procedure records in LPR2 were linked to
their admission or contact through the full record identifier; records with a
missing identifier were not linked. A data-quality audit on 4 October 2026
showed that an earlier linkage, which truncated long record identifiers, had
failed to link procedures for 355,581 (surgical procedures) and 1,576,198
(other procedures) distinct record identifiers. Correcting the linkage, and
dating procedures by the procedure start rather than the admission start, added
22,208 event labels and moved 5,484 event dates earlier, summed across seven
procedure outcomes (not distinct patients). The corresponding audit of LPR3
procedure-to-contact linkage found no lost records.

\paragraph{Data cutoff.} The administrative data cutoff was 19 March 2026 for
all endpoints, including second-line treatment; an earlier cutoff of 4 March
2026 was used only in historical sensitivity analyses. All endpoints shared the
same cutoff, which was checked when composite outcomes were formed. A death
recorded after the cutoff was treated as censoring at the cutoff, not as an
event.

% ===========================================================================
\section{Study population and index date (first-line treatment)}\label{sec:population}

The index date was the start of first-line treatment, reconstructed from the
treatment records by the following algorithm. Exploratory variants (disabling
anti-CD20 monotherapy as first-line treatment, and a fallback to
registry-recorded treatment dates) were not used.

\begin{enumerate}[label=\textbf{Step \arabic*.},leftmargin=4.2em]
\item \textbf{Diagnosis anchor.} For each patient, the earliest record of any
DALY-CARE haematological diagnosis (including MGUS, MBL and unspecified codes)
defined the diagnosis anchor date.
\item \textbf{Record filter.} Treatment records required a patient identifier,
a date and a drug code (or, if missing, a drug text). A record was retained if,
in order of precedence: (i) it originated from a quality registry (LYFO,
DaMyDa, DCLLR); (ii) its indication code started with a haematological ICD-10
prefix (DC81--DC89, DC90--DC96, DD47). It was discarded if (iii) its
four-character indication code was a listed non-haematological cancer (DC18,
DC19, DC20, DC25, DC34, DC43, DC50, DC56, DC61, DC71, DC78), or (iv) the
indication text matched a list of non-haematological Danish terms or the drug
text named a listed agent not used for lymphoid malignancy (e.g.\
fluorouracil, capecitabine, imatinib, azacitidine, hydroxycarbamide,
cisplatin, carboplatin). Otherwise (v) it was retained as of uncertain
indication.
\item \textbf{Search window.} Records dated from 30 days before to 3,650 days
after the diagnosis anchor were eligible.
\item \textbf{Episodes.} Records were grouped into one treatment episode while
successive treatment dates were at most 21 days apart (a gap rule, not a
maximum duration). The episode start was the earliest record date.
\item \textbf{Therapy-like criterion} (first matching rule): (a) the episode
contained an \emph{anchor} agent $\Rightarrow$ qualified; (b) it contained a
quality-registry record, unless it also contained an agent suspected of
prophylactic or non-lymphoma use and fewer than two distinct codes
$\Rightarrow$ qualified; (c) anti-CD20 monotherapy (ATC L01FA01, L01XC02, SKS
BOHJ11, or drug text rituximab or ofatumumab) on $\ge$2 dates
$\Rightarrow$ qualified; (d) it contained a suspect agent $\Rightarrow$
required $\ge$3 dates and $\ge$2 distinct codes; (e) otherwise $\ge$2 dates or
$\ge$2 distinct codes were required.

Anchor agents comprised the CLL agents (fludarabine, cyclophosphamide,
bendamustine, chlorambucil, obinutuzumab, ibrutinib, zanubrutinib,
acalabrutinib, idelalisib, venetoclax, alemtuzumab; ATC and SKS BWHA/BOHJ
codes), aggressive B-cell lymphoma agents and regimens (e.g.\ SKS codes for
CHOP, COP and CHOEP; mitoxantrone, ifosfamide, carboplatin, gemcitabine,
vincristine, melphalan, cytarabine, bortezomib, lenalidomide, thiotepa,
carmustine), myeloma regimens identified by drug text (e.g.\ daratumumab,
carfilzomib, pomalidomide, autologous stem-cell transplantation) and SKS
radiotherapy codes (codes described as ``strålebehandling'' [radiotherapy],
except listed prefixes, plus AFB04F3 and AFB03F3). Agents suspected of
prophylactic or non-lymphoma use were removed from the anchor list: rituximab,
unspecified CD20 antibody, methotrexate, gemcitabine and cisplatin identified
by text, doxorubicin, etoposide, oxaliplatin, selinexor, the generic cytostatic
codes BWHA1, BWHA2 and BWHA4, and bevacizumab.
\item \textbf{Curated CLL and MM anchors.} For patients in the curated CLL or
MM treatment-line data, the earliest first-line date (the earlier of the two if
present in both) was attached. An episode was an \emph{anchored episode} if
this date fell within it or within 30 days of its start or end. Anchored
episodes did not need to meet the therapy-like criterion.
\item \textbf{Selection.} Patients with a curated anchor received the anchored
episode closest to the curated date; all others received the earliest
qualifying episode. The index date was the start of the selected episode, not
the curated date. Patients without a qualifying episode had no index date and
were not included.
\end{enumerate}

\paragraph{Cohort flow.} Supplementary Table~\ref{tab:flow} shows the cohort
flow. Patients enter the outcome population only if they had a follow-up record
and both end of follow-up and the data cutoff fell after the index date; a
patient with a missing follow-up record, end of follow-up on or before the
index date (including death on the index date) or an index date on or after
the cutoff was therefore excluded from all outcomes. Twenty-six patients of the
first-line cohort were excluded in this way.

\begin{table}[ht]
\centering\footnotesize
\caption{Cohort flow.}\label{tab:flow}
\begin{tabular}{L{11cm}r}
\toprule
Stage & Patients\\
\midrule
Patients with a DALY-CARE haematological diagnosis (pretraining population) & 78,668\\
\quad No qualifying first-line treatment & 40,033\\
First-line index date identified (pass 1) & 38,635\\
\quad Disease assigned in pass 1 (diagnoses from 20 years before to 14 days after index) & 37,357\\
\quad Unassigned after pass 1 & 1,278\\
\quad Additionally assigned in pass 2 (coding-lag reassessment, up to 90 days after index) & 209\\
\quad Unassigned after pass 2 (excluded) & 1,069\\
Disease assigned after both passes & 37,566\\
\quad IgM MGUS excluded & 9\\
First-line cohort & 37,557\\
\quad No follow-up after the index date or index on/after the data cutoff & 26\\
Supervised population & 37,531\\
\quad Training (index date up to 31 December 2021) & 29,105\\
\quad Tuning (index date in 2022) & 1,980\\
\quad Held-out test (index date from 1 January 2023) & 6,446\\
\bottomrule
\end{tabular}
\end{table}

% ===========================================================================
\section{Disease classification}\label{sec:disease}

Diagnosis codes were mapped to disease groups using the DALY-CARE ICD-10
patterns with local refinements: SLL only by the exact code DC830;
lymphoplasmacytic lymphoma by DC830B or DC880; nodal marginal-zone lymphoma by
DC830C; splenic marginal-zone lymphoma by DC830D; the generic marginal-zone
group was not used; and DD472A (IgM MGUS) was mapped to an exclusion group.

\paragraph{Pass 1 (contemporaneous diagnoses).} Diagnoses dated from 7,300
days (about 20 years) before to 14 days after the index date were considered
(an asymmetric window). A disease group was \emph{confirmed} if any diagnosis
had a reliability rank of 3 or better (1 being the highest; quality-registry
diagnoses rank 1), if there were $\ge$2 diagnoses spanning $\ge$14 days, or if
there was a single rank-1 diagnosis. Among confirmed groups, unspecified and
supergroup-only labels (generic DALY-CARE, unspecified, unspecified B-cell
lymphoma, unspecified lymphoma and unspecified leukaemia) and MGUS/MBL were not
assignable. Two derived labels were added: Richter transformation (CLL/SLL and
DLBCL both confirmed in the window, with the earliest recorded DLBCL code
$\ge$90 days after the earliest recorded CLL/SLL code, both taken over the
patient's full history) and transformed follicular lymphoma (same rule for
FL). Specific labels took precedence over supergroups and aggressive over
non-aggressive labels; remaining ties were broken by aggressiveness (Richter;
then DLBCL, BL, PTCL, ALCL, AITL, LBL, PCL, MCL; then MM, solitary
plasmacytoma; then CLL, SLL, LPL; then FL, MZL, EMZL, HCL; then others), then
by best reliability rank, latest diagnosis date and specificity.

\paragraph{Pass 2 (coding-lag reassessment).} A specific lymphoma diagnosis is
often registered only after the diagnostic work-up is complete, which can be
after treatment has started. Patients with an index date but no assignable
disease after pass 1 therefore entered a second pass. Their earliest
confirmed, assignable diagnosis anywhere in their history became a new anchor;
first-line treatment was re-identified with the same algorithm in a window
from 90 days before this anchor to 20 years after it (curated CLL/MM anchors
still applied), and disease was re-assigned with a 10-year look-back and a
90-day forward window. Pass-2 patients therefore received a new index date,
replacing their pass-1 index. This allowed a specific diagnosis registered up
to 90 days after treatment start to define disease identity.

\paragraph{Regimen-based assignment and exclusions.} Patients still unassigned
after pass 2 could be assigned a disease group from their first-line regimen;
this added no patients to the final cohort. Patients assigned IgM MGUS (DD472A)
were excluded. Regimen-level characteristics (for example, salvage-like or
radiotherapy-only first episodes) were recorded descriptively and did not lead
to exclusion.

\paragraph{Disease taxonomy.} The analysis used 24 disease subtypes nested in
10 disease groups (Supplementary Table~\ref{tab:taxonomy}). The plasma-cell
dyscrasia supergroup was grouped with MM and the generic marginal-zone group
with indolent B-NHL; neither occurs among the 24 analysis subtypes.

\begin{table}[ht]
\centering\footnotesize
\caption{Disease groups and subtypes.}\label{tab:taxonomy}
\begin{tabular}{L{3.2cm}L{10.5cm}}
\toprule
Disease group & Subtypes\\
\midrule
CLL/SLL & CLL, SLL\\
Indolent B-NHL & FL, extranodal MZL, nodal MZL, splenic MZL, LPL, B-cell lymphoma without further specification\\
DLBCL-like & DLBCL, Richter transformation (coded or derived), transformed FL\\
MCL & MCL\\
HL & HL\\
MM & MM, plasma-cell leukaemia, solitary plasmacytoma\\
T-NHL & AITL, ALCL, PTCL, T-cell lymphoma without further specification\\
HCL & HCL\\
BL/LBL & BL, LBL\\
Amyloidosis & Amyloidosis\\
\bottomrule
\end{tabular}
\end{table}

% ===========================================================================
\section{Temporal split and prediction time}\label{sec:split}

\paragraph{Temporal split.} Patients were assigned by index date to
\textbf{training} (index date before 1 January 2022), \textbf{tuning} (1
January to 31 December 2022) or \textbf{held-out test} (from 1 January 2023).
Each patient had the same index date and split in every outcome; composite
outcomes required agreement of index date and split across components.

\paragraph{Prediction time.} The model input for each patient comprised all
records timestamped before 00:00 on the index date, followed by a prediction
token. Treatment and laboratory records on the index date itself were
therefore not model inputs, although an index-date laboratory draw counted
towards laboratory ascertainment (Section~\ref{sec:elig}). Because most
records carry a date but no time, events on the index date cannot be ordered
relative to treatment start; such events were not counted as outcomes, and
models were evaluated on events occurring at least 1 hour after the
prediction time.

\paragraph{Pretraining.} The foundation model was pretrained with an exclusive
calendar cutoff of 1 January 2022, applied in the same way, and its vocabulary
was derived from pre-2022 events of pretraining patients only. Pre-2022
histories of patients later assigned to tuning or held-out test were included
in pretraining; the evaluation therefore concerns future treatment episodes,
not unseen individuals.

\paragraph{Disease identity at prediction time.} Disease labels could rely on
codes registered up to 14 days (pass 1) or 90 days (pass 2) after the index
date. Models conditioned on disease therefore used information that was not
necessarily available at the index date for these patients; the main text
describes a sensitivity analysis restricted to labels available at treatment
start.

% ===========================================================================
\section{Disease-course endpoints}\label{sec:2l}

\subsection{Second-line treatment}
Second-line treatment was identified from the same filtered treatment records
as first-line treatment, with disease-specific rules.

\paragraph{Lymphoma groups} (DLBCL-like, HL, BL/LBL, MCL, T-NHL, HCL, indolent
B-NHL). (1) If the first-line episode spanned a single day, its end was
imputed as start plus an expected regimen duration. (2) Treatment records dated
more than a minimum gap after the first-line end were grouped into episodes
(21-day gap rule). (3) Episodes were excluded as autologous stem-cell
transplant consolidation (transplant or conditioning codes within the
consolidation window), consolidation (a share of first-line codes reappearing
at or above an overlap threshold within the consolidation window), maintenance
(only anti-CD20 or lenalidomide within the maintenance window), or single-date
anti-CD20. (4) Remaining episodes had to be therapy-like (as for first line,
except that quality-registry episodes needed $\ge$2 codes or $\ge$2 dates)
\emph{and} have at least moderate evidence of treatment for lymphoid
malignancy: quality-registry source; an anchor agent with haematological
indication; an anchor agent without suspect agents; or haematological
indication with $\ge$2 dates \emph{and} $\ge$2 distinct codes. Anti-CD20-only
episodes on $\ge$2 dates were therefore counted as retreatment and qualified
as second line when they met these evidence criteria (e.g.\ obinutuzumab, an
anchor agent, or haematological indication with two distinct codes), except
within the maintenance window in MCL and indolent B-NHL. The event date was
the start of the earliest qualifying episode. (5) A second-line date recorded
in LYFO after the index date replaced the derived date if it was earlier or if
no derived date existed. The disease-specific parameters are listed in
Supplementary Table~\ref{tab:2l}.

\begin{table}[ht]
\centering\footnotesize
\caption{Disease-specific parameters for second-line treatment in lymphoma groups.
Durations and windows in days; -- not applied.}\label{tab:2l}
\begin{tabular}{lrrrrrl}
\toprule
Group & Imputed first-line & Minimum & Consolidation & Maintenance & Overlap & Reliability\\
 & duration & gap & window & window & threshold & \\
\midrule
DLBCL-like & 126 & 60 & 90 & -- & 0.45 & high\\
HL & 120 & 60 & 90 & -- & 0.45 & high\\
BL/LBL & 28 & 45 & 90 & -- & 0.50 & low\\
MCL & 140 & 60 & 180 & 548 & 0.40 & moderate\\
T-NHL & 126 & 60 & 90 & -- & 0.40 & moderate\\
HCL & 7 & 90 & -- & -- & 0.30 & low\\
Indolent B-NHL & 168 & 90 & 120 & 730 & 0.35 & moderate\\
\bottomrule
\end{tabular}
\end{table}

\paragraph{Ascertainment in lymphoma groups.} Patients whose only
post-first-line retreatment evidence was weak were excluded from the risk set
rather than counted as event-free. Because second-line identification was
unreliable in BL/LBL and HCL, these patients were excluded unless a LYFO
second-line date created an event; BL/LBL, HCL and amyloidosis were
consequently excluded from all second-line-dependent outcomes.

\paragraph{CLL/SLL and MM.} The event was the earliest second-line date in
the curated CLL or MM treatment-line data (matched on disease group), if later
than the index date. CLL/SLL and MM patients without such a date were
event-free, assuming completeness of the curated data. Amyloidosis was not
included.

\paragraph{Censoring.} Second-line treatment was censored at the earlier of
the date of death and 19 March 2026; an event counted only if on or before
this date, and patients censored on or before the index date were excluded.
Death was not part of the endpoint and was handled as a competing event
(Section~\ref{sec:elig}). The endpoint is time to next qualifying treatment,
not directly observed progression.

\subsection{Treatment failure and transformation}
\textbf{Treatment failure} was the earlier of second-line treatment and death.
\textbf{Treatment failure or transformation} was the earliest of second-line
treatment, death and a new haematological malignancy. Transformation was
thus operationalised as a new three-character ICD-10 block among C81--C96 and
D45--D47 first recorded more than 31 days after the index date and not recorded
on or before index $+$31 days; it is not histologically confirmed
transformation and also captures unrelated second haematological neoplasms.
Both composites inherited the second-line ascertainment rules and the exclusion
of BL/LBL, HCL and amyloidosis (Section~\ref{sec:composites}).

\subsection{Overall survival and new malignancy}
Overall survival was the time to death on or before 19 March 2026, censored at
the administrative censoring date. A new haematological malignancy and a new
malignancy of any type were defined by the first post-index diagnosis in a
three-character ICD-10 block never recorded on or before index $+$31 days;
codes registered around treatment start, which typically reflect the presenting
disease, were thereby not counted. Relapse coded in the same block therefore
never counted, whereas a new block (for example, FL, C82, to DLBCL, C83) did.

{\footnotesize\setlength{\tabcolsep}{3pt}
\begin{longtable}{L{3.0cm}L{5.5cm}L{2.5cm}L{3.3cm}L{1.6cm}}
\caption{Disease-course endpoints (6 outcomes). OP, outcome population
(Section~\ref{sec:elig}).}\label{tab:course}\\
\toprule
Outcome & Definition & Source & Ascertainment and censoring & Competing event\\
\midrule\endfirsthead
\toprule
Outcome & Definition & Source & Ascertainment and censoring & Competing event\\
\midrule\endhead
\bottomrule\endfoot
Overall survival & Death on or before 19 March 2026. & Patient table & OP; administrative. & None\\
Second-line treatment & As above. & Treatment data, curated CLL/MM lines, LYFO & Disease-specific; censored at death or cutoff; excludes BL/LBL, HCL, amyloidosis. & Death\\
Treatment failure & Earliest of second-line treatment and death. & Composite & Intersection of components; earliest component censoring. & None\\
Treatment failure or transformation & Earliest of second-line treatment, death and new haematological malignancy. & Composite & As above. & None\\
New haematological malignancy & New ICD-10 block among C81--C96, D45, D46, D47 after index $+$31 days. & DALY-CARE haematological diagnoses & OP; administrative. & Death\\
New malignancy (any) & New ICD-10 block among C00--C97, except C44 and C80, after index $+$31 days. & DNPR diagnoses & OP; administrative. & Death\\
\end{longtable}}

% ===========================================================================
\section{Treatment-related and acute-care outcomes}\label{sec:outcomes}

\subsection{Outcome classes}
Each outcome belongs to one of the following classes.
\begin{description}[style=nextline,leftmargin=1.2em]
\item[Graded laboratory (DALY-CARE grading).] Laboratory values were graded by
the DALY-CARE standardised adverse-event grading algorithms, whose upstream
documentation describes several as derived from CTCAE v5. The event was the
first post-index measurement at or above the stated grade.
\item[Local laboratory threshold.] As above, with the numeric threshold stated
in the table.
\item[Baseline-relative laboratory.] The event was defined relative to the
patient's own pre-index value. The baseline was the median value on the latest
day with a measurement in the 365 days strictly before the index date; patients
without a gradable baseline were excluded. Many patients with lymphoid
malignancy have abnormal values before treatment (for example, through
hepatic or renal involvement); a fixed threshold would then record a
pre-existing abnormality at the first post-index draw as an event, whereas a
baseline-relative threshold requires worsening.
\item[Registry diagnosis or procedure.] The first post-index ICD-10 diagnosis
or SKS procedure code. For \emph{new-diagnosis} outcomes, a post-index code
counted only if its three-character block had never been recorded on or before
the index date (or index $+$ grace, where stated).
\item[Medication-use proxy.] Repeated antimicrobial administration or
dispensing (Section~\ref{sec:infections}).
\item[Composite.] Earliest of its components (Section~\ref{sec:composites}).
\end{description}

Except for the baseline-relative definitions and the alkaline phosphatase
tiers, laboratory events are the first post-index documentation of an abnormal
value, not new onset or worsening: abnormal baseline values were not excluded.
No outcome is a verified CTCAE v6 grade; grade labels denote the output of the
grading algorithm (``CTCAE-aligned''). Prediction horizons were 30, 60, 90, 180,
365 and 730 days (primary 365 days). For every outcome in
Supplementary Tables~\ref{tab:inf}--\ref{tab:acute}, death was a competing event.

\paragraph{Abbreviations in Supplementary Tables~\ref{tab:inf}--\ref{tab:acute}.}
\emph{Sources}: DNPR-Dx, National Patient Registry diagnoses; DNPR-Proc,
National Patient Registry procedures (LPR2 and LPR3); Lab, laboratory results;
SMR, Hospital Medicines Register; Rx, National Prescription Registry.
\emph{Baseline}: Lab rule, shared laboratory ascertainment rule
(Section~\ref{sec:labrule}); OP, all patients in the outcome population.
\emph{Ascertainment}: Admin, censored at the administrative censoring date;
Lab, laboratory censoring (Section~\ref{sec:labrule}).
\emph{Class}: Grd, DALY-CARE grading; Loc, local threshold; Rel,
baseline-relative; Reg, registry; Prx, medication-use proxy; Cmp, composite.

\subsection{Infection and antimicrobial use}\label{sec:infections}
Diagnosis-based infection outcomes used the first post-index code without
exclusion of prior episodes, because infections are acute events and a new
post-index episode is clinically relevant even after an earlier one. The
medication-use proxies grouped administrations (SMR, intravenous route only)
or dispensings (prescription registry) of antimicrobials (DALY-CARE ATC
selection) into episodes separated by gaps of more than 7 days; an episode
qualified if it spanned at least 1 day, that is, $\ge$2 administrations or
dispensings on different days. The event was the first date of the first
qualifying post-index episode. No prophylaxis exclusion was applied to
intravenous administrations; for prescriptions, ATC J01EE01
(sulfamethoxazole--trimethoprim), J05AB01 (aciclovir), J05AB11
(valaciclovir) and the exact code P02 were excluded as typical prophylaxis.
These outcomes measure repeated antimicrobial use, not confirmed infection,
days of supply or treatment intent. Because SMR reporting began on 15 May 2018,
absence of intravenous administrations before that date is not evidence that
none occurred; the intravenous proxy and the composite that contains it were
therefore restricted to patients with index date on or after 15 May 2018.

{\footnotesize\setlength{\tabcolsep}{3pt}\otablebegin
\caption{Infection and antimicrobial-use outcomes (6 outcomes).}\label{tab:inf}\\
\otablehead
Sepsis & ICD-10 A40, A41; first post-index code. & DNPR-Dx & OP & Admin & Reg\\
Serious infection & ICD-10 A40, A41, B44, B45, B46, B59 (sepsis, invasive fungal infection, pneumocystosis) as one code set. & DNPR-Dx & OP & Admin & Reg\\
Antibiotic treatment (procedure) & SKS procedure codes beginning BPH; dated by procedure start. & DNPR-Proc & OP & Admin & Reg\\
Intravenous antimicrobial use & Qualifying episode of intravenous antimicrobial administration. & SMR & OP; index date $\ge$ 15 May 2018 & Admin & Prx\\
Prescription antimicrobial use & Qualifying episode of antimicrobial dispensing, excluding prophylaxis codes. & Rx & OP & Admin & Prx\\
Infection diagnosis or antimicrobial use & Earliest of serious infection, prescription and intravenous antimicrobial use. & Composite & Intersection; index date $\ge$ 15 May 2018 & Earliest component & Cmp\\
\end{longtable}}

\subsection{Blood counts, cytopenia and transfusion}
{\footnotesize\setlength{\tabcolsep}{3pt}\otablebegin
\caption{Blood count, cytopenia and transfusion outcomes (16 outcomes).}\label{tab:blood}\\
\otablehead
Anaemia, grade $\ge$2; grade $\ge$3 & Haemoglobin graded $\ge$2 / $\ge$3. & Lab & Lab rule & Lab & Grd\\
Neutropenia, grade $\ge$2; grade $\ge$3 & Neutrophil count decreased, $\ge$2 / $\ge$3. & Lab & Lab rule & Lab & Grd\\
Leukopenia, grade $\ge$2; grade $\ge$3 & White blood cell count decreased, $\ge$2 / $\ge$3. & Lab & Lab rule & Lab & Grd\\
Thrombocytopenia, grade $\ge$2; grade $\ge$3 & Platelet count decreased, $\ge$2 / $\ge$3 (grade 3/4 boundary 25$\times10^9$/L, CTCAE v5). & Lab & Lab rule & Lab & Grd\\
Lymphocytosis, grade $\ge$2; grade $\ge$3 & Lymphocyte count increased, $\ge$2 / $\ge$3. In CLL largely reflects disease activity. & Lab & Lab rule & Lab & Grd\\
Eosinophilia & Eosinophil count increased, grade $\ge$1. & Lab & Lab rule & Lab & Grd\\
Haptoglobin decreased & Haptoglobin decreased (binary grading). Not a diagnosis of haemolysis. & Lab & Lab rule & Lab & Grd\\
Severe cytopenia & Earliest of neutropenia, thrombocytopenia and anaemia, each grade $\ge$3. & Composite & Intersection & Earliest component & Cmp\\
Red-cell transfusion & SKS BOQA0, BOQA00--BOQA05; dated by procedure start. & DNPR-Proc & OP & Admin & Reg\\
Platelet transfusion & SKS BOQC, BOQC3. & DNPR-Proc & OP & Admin & Reg\\
Any transfusion & Union of the red-cell and platelet code sets (single code-set outcome). & DNPR-Proc & OP & Admin & Reg\\
\end{longtable}}

\subsection{Renal}
{\footnotesize\setlength{\tabcolsep}{3pt}\otablebegin
\caption{Renal outcomes (5 outcomes).}\label{tab:renal}\\
\otablehead
Acute kidney injury (first documented) & First post-index start of an acute kidney injury episode derived from creatinine by the DALY-CARE algorithm. & Lab & Lab rule & Lab & Grd\\
Creatinine increased, grade $\ge$2; grade $\ge$3 & Creatinine graded $\ge$2 / $\ge$3. & Lab & Lab rule & Lab & Grd\\
Creatinine rise $>$1.5$\times$ pre-index baseline & Creatinine (\textmu mol/L) $>$1.5 times baseline on any post-index day. Not a CTCAE grade or an acute kidney injury diagnosis. & Lab & Lab rule; baseline value $>$0 (no reference-limit requirement) & Lab & Rel\\
Dialysis & SKS procedure codes beginning BJFD. & DNPR-Proc & OP & Admin & Reg\\
\end{longtable}}

\subsection{Hepatic and coagulation}
{\footnotesize\setlength{\tabcolsep}{3pt}\otablebegin
\caption{Hepatic and coagulation outcomes (9 outcomes). ULN, upper limit of normal.}\label{tab:liver}\\
\otablehead
ALT increased, grade $\ge$2; grade $\ge$3 & ALT (U/L), with reference ULN $>$0. If baseline $>$ its ULN: value $>$1.5 / $>$2.0 times baseline. Otherwise: value $>$3 / $>$5 times the ULN of the measurement. & Lab & Lab rule; gradable baseline (value $>$0, baseline ULN $>$0) & Lab & Rel\\
Alkaline phosphatase elevation, tier $\ge$2; tier $\ge$3 & $>$2.5 / $>$5 times the sample ULN if the running mean of prior samples was normal, otherwise $>$2.5 / $>$5 times that prior mean; missing ULN set to 105 U/L. Time-varying baseline; CTCAE v6 grades ALP at grade 1 only. & Lab & Lab rule & Lab & Grd\\
Bilirubin increased, grade $\ge$2; grade $\ge$3 & Bilirubin graded $\ge$2 / $\ge$3. & Lab & Lab rule & Lab & Grd\\
Hepatotoxicity & Earliest of ALT increased grade $\ge$3 (baseline-relative), AST increased grade $\ge$3 and bilirubin increased grade $\ge$3. & Composite & Intersection, hence gradable ALT baseline & Earliest component & Cmp\\
INR increased, grade $\ge$2; grade $\ge$3 & INR graded $\ge$2 / $\ge$3; anticoagulant therapy not considered. & Lab & Lab rule & Lab & Grd\\
\end{longtable}}

\subsection{Electrolytes and metabolic}
{\footnotesize\setlength{\tabcolsep}{3pt}\otablebegin
\caption{Electrolyte and metabolic outcomes (19 outcomes). LLN/ULN, lower/upper limit of normal of the measurement.}\label{tab:metab}\\
\otablehead
Hyponatraemia, grade $\ge$2 & Sodium $<$130 mmol/L and $<$ LLN (LLN 136 mmol/L if missing). Symptoms not observed. & Lab & Lab rule & Lab & Loc\\
Sodium $<$125 mmol/L & Sodium $<$125 mmol/L and $<$ LLN. Does not capture symptomatic hyponatraemia at 125 to $<$130 mmol/L. & Lab & Lab rule & Lab & Loc\\
Potassium $<$3.0 mmol/L & Potassium $<$3.0 mmol/L on any post-index day. Symptoms and intervention not observed. & Lab & Lab rule & Lab & Loc\\
Hypernatraemia, grade $\ge$2; grade $\ge$3 & Sodium graded $\ge$2 / $\ge$3. & Lab & Lab rule & Lab & Grd\\
Hyperkalaemia, grade $\ge$2; grade $\ge$3 & Potassium graded $\ge$2 / $\ge$3. & Lab & Lab rule & Lab & Grd\\
Hypercalcaemia, grade $\ge$2; grade $\ge$3 & Calcium graded $\ge$2 / $\ge$3. & Lab & Lab rule & Lab & Grd\\
Hypocalcaemia, grade $\ge$2; grade $\ge$3 & Calcium graded $\ge$2 / $\ge$3 (albumin correction not established). & Lab & Lab rule & Lab & Grd\\
Hypomagnesaemia, grade $\ge$2 & Magnesium (NPU02647) $<$0.5 mmol/L and $<$ LLN (LLN 0.7 mmol/L if missing). & Lab & Lab rule & Lab & Loc\\
Hypertriglycer\-idaemia, grade $\ge$2; grade $\ge$3 & Triglycerides graded $\ge$2 / $\ge$3; fasting status unknown. & Lab & Lab rule & Lab & Grd\\
Hypoalbuminaemia, grade $\ge$2; grade $\ge$3 & Albumin graded $\ge$2 / $\ge$3. & Lab & Lab rule & Lab & Grd\\
Hyperphosphataemia & Phosphate $>$ ULN (ULN 1.45 mmol/L if missing); NPU03096 used where the standard phosphate code was absent. Laboratory elevation only. & Lab & Lab rule & Lab & Loc\\
Hyperuricaemia & Urate $>$ ULN; no event if the ULN was missing. Laboratory elevation only. & Lab & Lab rule & Lab & Loc\\
LDH increased & Lactate dehydrogenase $>$ ULN; no event if the ULN was missing. & Lab & Lab rule & Lab & Loc\\
\end{longtable}}

\subsection{Cardiovascular and thrombotic}
New-diagnosis rule without grace period: a post-index code counted only if its
three-character block had not been recorded on or before the index date. Blocks
were evaluated separately, so that, for example, a prior I80 did not exclude a
new I26.

{\footnotesize\setlength{\tabcolsep}{3pt}\otablebegin
\caption{Cardiovascular and thrombotic outcomes (4 outcomes).}\label{tab:cv}\\
\otablehead
Atrial fibrillation & ICD-10 I48; new diagnosis. & DNPR-Dx & OP & Admin & Reg\\
Heart failure & ICD-10 I50; new diagnosis. & DNPR-Dx & OP & Admin & Reg\\
Stroke or TIA & ICD-10 I63, I64, G45; new diagnosis per block. & DNPR-Dx & OP & Admin & Reg\\
Venous thromboembolism & ICD-10 I26, I80, I81, I82; new diagnosis per block. & DNPR-Dx & OP & Admin & Reg\\
\end{longtable}}

\subsection{Acute care and other events}
\paragraph{Acute hospitalisation.} The definition followed the Danish database
of acute hospital contacts. In LPR2, inpatient admissions (acute or planned)
were combined with outpatient and emergency-room contacts with an acute mode
of admission (LPR2 patient type: inpatient, outpatient, or emergency room
before 2014; from 2014 emergency-room contacts were registered as outpatient
contacts with acute mode of admission). Missing admission and discharge hours
were set to 09:00 and 21:00; outpatient and emergency contacts were treated as
point events. In LPR3, physical attendances were used, acute if registered
with acute priority; planned LPR3 contacts were retained only if they lasted
$\ge$12 hours. Psychiatric specialties were excluded in both eras. Contacts
separated by at most 4 hours were merged into one episode. An episode was
\emph{acute} if any of its contacts was acute, and a \emph{hospitalisation} if
it contained an LPR2 inpatient admission, lasted $\ge$12 hours or crossed
midnight. The outcome comprised acute hospitalisations only, so a planned
admission counted only if it was merged with an acute contact; planned
admissions, which typically deliver treatment, were not outcome events. The
event date was the start date of the first post-index acute hospitalisation.

{\footnotesize\setlength{\tabcolsep}{3pt}\otablebegin
\caption{Acute-care and other outcomes (4 outcomes).}\label{tab:acute}\\
\otablehead
Acute hospitalisation & As described above. & DNPR admissions and contacts & OP & Admin & Reg\\
Emergency department stay & SKS procedure codes beginning NAB, excluding codes containing NAB1. & DNPR-Proc & OP & Admin & Reg\\
Respiratory support & SKS BGDA0, BGDA1, BGDA2, BGDA7. & DNPR-Proc & OP & Admin & Reg\\
Bone fracture & ICD-10 S02, S12, S22, S32, S42, S52, S62, S72, S82, S92, T02, M80, M84.3, M84.4 (including subcodes); first post-index code, no exclusion of prior fractures. & DNPR-Dx & OP & Admin & Reg\\
\end{longtable}}

% ===========================================================================
\section{Outcome-specific eligibility, ascertainment and censoring}\label{sec:elig}

\paragraph{Outcome population.} All outcomes started from the first-line
cohort restricted to patients with a follow-up record whose end of follow-up
and data cutoff both fell after the index date (37,531 patients;
Section~\ref{sec:population}).

\subsection{Laboratory ascertainment}\label{sec:labrule}
Laboratory outcomes are only observable in patients whose blood is tested.
Absence of an abnormal value in a patient without measurements is not evidence
that the abnormality did not occur. All graded, local-threshold and
baseline-relative laboratory outcomes and acute kidney injury therefore used
a shared ascertainment rule based on the dates of \emph{any} laboratory
measurement (not only the outcome analyte). A patient was eligible if all three
conditions held:
\begin{itemize}
\item \emph{source coverage}: index date on or after 1 January 2017 (the date
from which laboratory capture was regarded as complete), or any measurement
between 365 days before the index date and the administrative censoring date;
\item \emph{baseline}: any measurement between 365 days before the index date
and the index date inclusive;
\item \emph{entry}: any measurement between the index date and 30 days after
it, inclusive, \emph{or} death within 30 days of the index date with index
date on or after 1 January 2017 (so that early death did not remove patients
from the risk set).
\end{itemize}
A single measurement on the index date satisfied both the baseline and the
entry condition. Observing an event did not override these criteria.

\paragraph{Laboratory censoring.} Patients with an event were censored at the
administrative censoring date. Event-free patients were censored at the
earlier of their last post-index measurement of any analyte plus 180 days and
the administrative censoring date, so that time without laboratory monitoring
was not credited as event-free. The ascertainment status of each patient was
recorded (before laboratory coverage, missing baseline, no post-index
measurement, insufficient monitoring, eligible with event, eligible with
competing death, eligible event-free).

\paragraph{Baseline-relative and supplementary threshold outcomes.} The risk set was that of
the corresponding graded outcome (creatinine increased grade $\ge$2; ALT
increased grade $\ge$2; for potassium $<$3.0 mmol/L, hypokalaemia grade
$\ge$3), further restricted to patients with a gradable pre-index baseline
where required, and censored by the laboratory rule.

\subsection{Non-laboratory outcomes}
Diagnosis, procedure, hospitalisation, medication-use and survival outcomes
were administratively ascertained: all patients in the outcome population were
eligible and censored at the administrative censoring date. The only calendar
restriction was the SMR start date (index date on or after 15 May 2018) for
the intravenous antimicrobial-use proxy and the infection composite that
contains it. A previous restriction of procedure outcomes to the LPR3 era (go-live 3
February 2019 plus 30 days) was removed after the procedure linkage and dating
were corrected (Section~\ref{sec:sources}); hospitalisation and diagnosis
outcomes were deliberately not restricted by era (decision of 5 October 2026).
Calendar trends in registry capture are therefore part of these labels. For
every outcome, the released label set itself defined the risk set; no
additional eligibility masks were applied.

\subsection{Competing death and horizon labels}
For a fixed horizon $h$, a patient was a case if the event occurred by $h$ and
before any competing death; a known non-case if follow-up reached $h$ or a
competing death occurred by $h$; and otherwise of unknown status
(right-censored) and excluded from the fixed-horizon loss and from binary
metrics. Death precludes the subsequent occurrence of every non-fatal outcome;
treating it as ordinary censoring would overstate the risk of these outcomes,
whereas treating it as a known non-case keeps patients who die early in the
denominator. A death counted as a competing event only if it occurred on or
after the index date, on or before the outcome's censoring date and within
the horizon. For laboratory outcomes, an event-free patient who died more than
180 days after the last measurement was therefore censored rather than
counted as a competing death. Endpoints that include death (overall survival
and both treatment-failure composites) had no competing event.

% ===========================================================================
\section{Selection of the outcome panel}\label{sec:panel}

The panel was selected from a historical list of 87 candidate outcomes using
only training and tuning data. Support was the number of eligible patients
with an event within 365 days of the index date, without a follow-up
requirement. The minimum support was $\ge$100 events in training and $\ge$20
in tuning. Held-out test counts were reported descriptively and never used for
selection. Clinical decisions took precedence over support
(Supplementary Table~\ref{tab:panel}); every released outcome was re-checked
against the support threshold.

{\footnotesize
\setlength{\tabcolsep}{3pt}
\begin{longtable}{L{3.6cm}rL{12.0cm}}
\caption{Selection of the outcome panel from 87 to 69 outcomes.}\label{tab:panel}\\
\toprule
Step & $n$ & Outcomes and reason\\
\midrule\endfirsthead
\toprule
Step & $n$ & Outcomes and reason\\
\midrule\endhead
Historical candidate list & 87 & \\
Retired & $-$6 & Hypofibrinogenaemia grade $\ge$2 and grade $\ge$3 (sparse; assay, unit and baseline grading not verified); leukocytosis grade $\ge$2 (the grading algorithm requested grade 3 despite the label, and CTCAE v6 discourages the term for leukaemia-related rises); DIC-like coagulopathy, Hy's-law drug-induced liver injury and tumour lysis syndrome (laboratory co-occurrence proxies without the usual monitoring-based ascertainment).\\
Definition concerns, excluded & $-$5 & Lymphopenia grade $\ge$2 and grade $\ge$3 (CTCAE v6 has no grade 2/3 tier); hypokalaemia grade $\ge$2 and grade $\ge$3 (local thresholds $\le$3.0 and $\le$2.5 mmol/L, not CTCAE v6); major adverse cardiovascular events (composite included atrial fibrillation, venous thromboembolism and heart failure).\\
Below support threshold & $-$9 & AST increased grade $\ge$2 and grade $\ge$3; herpes zoster; hypermagnesaemia grade $\ge$3; hypoglycaemia grade $\ge$2 and grade $\ge$3; hypomagnesaemia grade $\ge$3; invasive fungal infection; myocardial infarction.\\
Primary candidates & 60 & Met the support threshold; no clinical concern.\\
Definition concerns, retained & $+$7 & Renamed to describe the measured quantity: alkaline phosphatase grade $\ge$2/$\ge$3 as alkaline phosphatase elevation tier $\ge$2/$\ge$3; hyponatraemia grade $\ge$3 as sodium $<$125 mmol/L; intravenous, prescription and combined infection outcomes as antimicrobial-use proxies. Hyponatraemia grade $\ge$2 kept its name. All met the support threshold.\\
Replacements & 0 & ALT increased grade $\ge$2/$\ge$3 replaced by the baseline-relative definition; hepatotoxicity rebuilt with the baseline-relative ALT grade $\ge$3.\\
Additions & $+$2 & Creatinine rise $>$1.5$\times$ pre-index baseline; potassium $<$3.0 mmol/L.\\
\midrule
Final panel & 69 & \\
\bottomrule
\end{longtable}}

Candidate tumour lysis labels (shared laboratory ascertainment; $\ge$2 of urate
$\ge$476\,\textmu mol/L, potassium $\ge$6.0 mmol/L, phosphate $\ge$1.45 mmol/L
and calcium $\le$1.75 mmol/L on the same day, on days 1--7 after the index
date; with a variant additionally involving kidney injury) remained
exploratory and outside the panel, with 243/45/114 and 207/41/100 events
(training/tuning/held-out test), respectively.

% ===========================================================================
\section{Composite outcomes, endpoint groups and display groups}\label{sec:composites}

\paragraph{Composite rule.} A patient was eligible for a composite only if
eligible for every component (intersection of risk sets). The composite was
censored at the earliest component censoring date, and the event was the
earliest component event, counted only if on or before that censoring date.
Components had to agree on index date and split. Death was an event only when
overall survival was a component; otherwise death remained censoring in each
component and was handled as a competing event. The composites were treatment
failure, treatment failure or transformation, severe cytopenia, hepatotoxicity,
and infection diagnosis or antimicrobial use. Any transfusion and serious
infection were defined as single code-set outcomes rather than composites.

\paragraph{Endpoint groups.} Training loss and checkpoint selection used a
macro-average over endpoint groups of the competing-risk, cumulative-incidence
based area under the curve on the tuning set. Endpoint groups combined the
grade variants of one endpoint (grade $\ge$2 and grade $\ge$3), the two
alkaline phosphatase tiers, and hyponatraemia grade $\ge$2 with sodium
$<$125 mmol/L, giving 52 endpoint groups for the 69 outcomes. Related
diagnoses, proxies and composites were not merged.

\paragraph{Display groups.} For figures only, the 69 outcomes were assigned
once each to 17 descriptive groups; creatinine rise $>$1.5$\times$ baseline
was placed under ``Kidney injury and renal support'' and potassium
$<$3.0 mmol/L under ``Sodium and potassium''. Display groups did not affect
model weighting.

% ===========================================================================
\section{Limitations of the outcome definitions}\label{sec:limits}
\begin{enumerate}
\item \textbf{Proxies rather than adjudicated events.} Antimicrobial-use
outcomes measure repeated antimicrobial administration or dispensing, not
confirmed infection; the prescription proxy has no calendar restriction and
intravenous prophylaxis was not excluded. Transformation is a new
haematological ICD-10 block, not histologically confirmed transformation.
Second-line treatment is time to next qualifying treatment, not observed
progression.
\item \textbf{CTCAE-aligned, not CTCAE.} Laboratory grades derive from the
DALY-CARE grading algorithms (several described as CTCAE v5-derived) and local
thresholds, some using reference limits with fixed fallback values. Symptoms,
interventions and clinical context required by CTCAE for several grades are
not observed; no outcome is a verified CTCAE v6 grade.
\item \textbf{Prevalent versus incident abnormalities.} Except for the
baseline-relative definitions, laboratory events are the first post-index
documentation of an abnormal value, not incident toxicity.
\item \textbf{Analyte-agnostic laboratory ascertainment.} Eligibility and
censoring use measurements of any analyte, so a patient can be eligible for an
outcome whose analyte was rarely measured.
\item \textbf{Coding lag and retrospective disease identity.} Disease labels
can use codes registered up to 14 days (pass 1) or 90 days (pass 2) after the
index date and are used as model input.
\item \textbf{First-line identification.} The search is anchored at the first
DALY-CARE diagnosis of any kind and limited to 10 years in pass 1; treatment
records without indication are retained as uncertain; curated anchors exist
for CLL and MM only.
\item \textbf{Registry era.} Diagnosis, procedure and hospitalisation outcomes
are not restricted across LPR2 and LPR3, so changes in coding practice over
calendar time are part of the labels; for LPR3 contacts without recorded clock
times, the 12-hour criterion reduces to ending on a later calendar day.
Diagnosis-based infection and fracture outcomes do not exclude continued coding
of a pre-index episode.
\item \textbf{Second-line ascertainment.} Second-line treatment depends on the
completeness of the curated CLL and MM data, excludes patients with only weak
retreatment evidence (enriching the risk set for events) and is censored at
death or the data cutoff rather than at end of follow-up as for other
outcomes.
\item \textbf{Tumour lysis syndrome} is not part of the panel; candidate
laboratory definitions remain exploratory and the historical co-occurrence
proxy was retired.
\end{enumerate}

\bibliographystyle{unsrtnat}
\bibliography{paper}
\end{document}
